\subsection{Modal Characteristics}
Table~\ref{tab:modal_period_mass_z040} compares the first three natural periods and the dominant translational modal mass participation of the seven building configurations under Soil Type C at $Z=0.40$. This common demand case permits comparison of their dynamic characteristics.

\begin{table}[!htbp]
\centering
\caption{First three modal periods and dominant translational mass participation for Soil Type C, $Z=0.40$.}
\label{tab:modal_period_mass_z040}
\small
\begin{tabular}{@{}lccccc@{}}
\toprule
Model & $T_1$ (s) & $T_2$ (s) & $T_3$ (s) & Mode 1 UX (\%) & Mode 2 UY (\%) \\
\midrule
Regular & 0.699 & 0.661 & 0.569 & 80.32 & 80.12 \\
L & 0.698 & 0.660 & 0.564 & 79.95 & 80.42 \\
L1 & 0.685 & 0.649 & 0.561 & 79.14 & 80.16 \\
L2 & 0.649 & 0.593 & 0.476 & 79.83 & 79.07 \\
T & 0.698 & 0.674 & 0.569 & 79.09 & 74.82 \\
U & 0.707 & 0.649 & 0.569 & 79.30 & 78.77 \\
Non-Parallel & 0.650 & 0.599 & 0.489 & 79.14 & 78.75 \\
\bottomrule
\end{tabular}
\end{table}


The first mode is predominantly X-translational for every configuration, contributing 79.09--80.32\% of the translational mass in UX. The second mode is predominantly Y-translational, contributing 74.82--80.42\% in UY. The fundamental period $T_1$ ranges from 0.649~s for L2 to 0.707~s for U; $T_2$ ranges from 0.593 to 0.674~s, while $T_3$ ranges from 0.476 to 0.569~s. These differences indicate modest but measurable changes in the dynamic characteristics among the investigated plan configurations and do not by themselves establish a performance ranking.

\begin{table}[!htbp]
\centering
\caption{Cumulative translational modal mass participation for Soil Type C, $Z=0.40$.}
\label{tab:modal_cumulative_mass_z040}
\small
\resizebox{\textwidth}{!}{%
\begin{tabular}{@{}lcccccc@{}}
\toprule
Model & SumUX at Mode 3 (\%) & SumUY at Mode 3 (\%) & X $\geq$ 90\% mode & SumUX (\%) & Y $\geq$ 90\% mode & SumUY (\%) \\
\midrule
Regular & 80.67 & 81.86 & 4 & 92.39 & 5 & 91.89 \\
L & 80.18 & 81.35 & 4 & 92.25 & 5 & 91.96 \\
L1 & 79.38 & 80.56 & 4 & 91.94 & 5 & 91.71 \\
L2 & 79.84 & 81.49 & 4 & 91.96 & 5 & 92.03 \\
T & 79.51 & 80.75 & 4 & 92.04 & 5 & 90.94 \\
U & 79.94 & 81.60 & 4 & 92.01 & 5 & 91.38 \\
Non-Parallel & 80.44 & 81.90 & 4 & 92.31 & 5 & 92.39 \\
\bottomrule
\end{tabular}
}
\end{table}


For every configuration, cumulative translational participation exceeds 90\% by Mode~4 in the X direction and by Mode~5 in the Y direction. The corresponding cumulative values at the threshold modes are approximately 91.94--92.39\% in X and 90.94--92.39\% in Y. By Mode~12, cumulative translational mass participation reaches 100\% in both directions for all seven configurations. The modal results therefore provide an explicit direction-wise participation check in addition to the period comparison used elsewhere in the response discussion.
\FloatBarrier
